\documentclass[UTF8,fontset=fandol,aspectratio=169,11pt]{ctexbeamer}
% Shared Beamer preamble for Big Data and Management Decision-Making slides.

\usetheme{Madrid}
\usecolortheme{default}
\usefonttheme{professionalfonts}

\usepackage{amsmath,amssymb,mathtools}
\usepackage{booktabs,array}
\usepackage{tikz}
\usepackage{graphicx}
\usepackage{listings}
\usepackage{xcolor}
\usetikzlibrary{arrows.meta,positioning,shapes.geometric}

\newcommand{\BDMFandolPath}{/usr/local/texlive/2025/texmf-dist/fonts/opentype/public/fandol/}
\setCJKmainfont[
  Path=\BDMFandolPath,
  UprightFont=FandolSong-Regular.otf,
  BoldFont=FandolSong-Regular.otf,
  ItalicFont=FandolKai-Regular.otf,
  BoldItalicFont=FandolKai-Regular.otf
]{FandolSong-Regular.otf}
\setCJKsansfont[
  Path=\BDMFandolPath,
  UprightFont=FandolSong-Regular.otf,
  BoldFont=FandolSong-Regular.otf
]{FandolSong-Regular.otf}
\setCJKmonofont[
  Path=\BDMFandolPath,
  UprightFont=FandolSong-Regular.otf,
  BoldFont=FandolSong-Regular.otf
]{FandolSong-Regular.otf}
\setmonofont{Menlo}

\definecolor{AcademicBlue}{RGB}{22,44,73}
\definecolor{AcademicTeal}{RGB}{22,119,119}
\definecolor{AcademicGray}{RGB}{76,84,95}
\definecolor{AcademicLight}{RGB}{245,247,250}
\definecolor{AcademicAmber}{RGB}{181,111,35}
\definecolor{CodeBg}{RGB}{248,248,248}

\setbeamercolor{structure}{fg=AcademicBlue}
\setbeamercolor{frametitle}{fg=white,bg=AcademicBlue}
\setbeamercolor{title}{fg=white,bg=AcademicBlue}
\setbeamercolor{block title}{fg=white,bg=AcademicBlue}
\setbeamercolor{block body}{bg=AcademicLight,fg=black}
\setbeamertemplate{navigation symbols}{}
\setbeamertemplate{footline}[frame number]
\setbeamertemplate{itemize item}{\raisebox{0.2ex}{\(\blacktriangleright\)}}
\setbeamerfont{title}{series=\mdseries}
\setbeamerfont{subtitle}{series=\mdseries}
\setbeamerfont{frametitle}{series=\mdseries}
\setbeamerfont{block title}{series=\mdseries}
\setbeamerfont{section in toc}{series=\mdseries}

\lstdefinestyle{pythonstyle}{
  basicstyle=\ttfamily\scriptsize,
  backgroundcolor=\color{CodeBg},
  keywordstyle=\color{AcademicBlue},
  commentstyle=\color{AcademicGray},
  stringstyle=\color{AcademicAmber},
  showstringspaces=false,
  breaklines=true,
  frame=single,
  rulecolor=\color{black!20},
  columns=fullflexible
}

\lstdefinestyle{sqlstyle}{
  language=SQL,
  basicstyle=\ttfamily\scriptsize,
  backgroundcolor=\color{CodeBg},
  keywordstyle=\color{AcademicBlue},
  commentstyle=\color{AcademicGray},
  stringstyle=\color{AcademicAmber},
  showstringspaces=false,
  breaklines=true,
  frame=single,
  rulecolor=\color{black!20},
  columns=fullflexible
}

\lstset{style=pythonstyle}

\hypersetup{
  colorlinks=true,
  linkcolor=AcademicBlue,
  urlcolor=AcademicTeal,
  pdfauthor={大数据与管理决策基础},
  pdfsubject={大数据与管理决策基础课程课件}
}


\title[SQL、KPI 与指标层]{第3讲：SQL、KPI 与可审计指标层}
\subtitle{连接、聚合、窗口函数与管理指标口径}
\author{《大数据与管理决策基础》}
\date{}

\lstset{style=sqlstyle}

\begin{document}

\begin{frame}
  \titlepage
  \begin{center}
    \small 核心命题：SQL 不是取数脚本，而是管理指标口径的可执行表达。
  \end{center}
\end{frame}

\begin{frame}{本讲学习目标}
\begin{enumerate}
  \item 理解 SQL 查询与关系模型的关系；
  \item 用连接、过滤和聚合计算经营 KPI；
  \item 区分事实粒度、分析粒度和展示粒度；
  \item 正确处理分子、分母和比率型指标；
  \item 使用窗口函数表达排名、累计和移动窗口；
  \item 将 SQL 上升为可治理的指标层。
\end{enumerate}
\end{frame}

\begin{frame}{课程主线中的第3周}
\begin{center}
\resizebox{0.95\linewidth}{!}{%
\begin{tikzpicture}[
  node/.style={draw=AcademicBlue!65, fill=AcademicLight, rounded corners=1pt,
               minimum height=0.8cm, minimum width=2.2cm, align=center},
  arrow/.style={-{Stealth[length=2mm]}, thick, AcademicTeal}
]
\node[node] (obj) {业务对象};
\node[node, right=0.9cm of obj] (grain) {事实粒度};
\node[node, right=0.9cm of grain] (sql) {SQL 查询};
\node[node, right=0.9cm of sql] (kpi) {指标口径};
\node[node, right=0.9cm of kpi] (act) {行动审计};
\draw[arrow] (obj) -- (grain);
\draw[arrow] (grain) -- (sql);
\draw[arrow] (sql) -- (kpi);
\draw[arrow] (kpi) -- (act);
\end{tikzpicture}
}
\end{center}

\begin{alertblock}{判断标准}
查询能运行，只说明语法成立；指标能被复核，才说明管理口径成立。
\end{alertblock}
\end{frame}

\begin{frame}[fragile]{一个区域收入查询}
\begin{lstlisting}
SELECT
  c.region,
  SUM(o.quantity * p.unit_price) AS revenue
FROM orders AS o
JOIN customers AS c
  ON o.customer_id = c.customer_id
JOIN products AS p
  ON o.product_id = p.product_id
WHERE o.order_date >= DATE '2026-03-01'
  AND o.order_date <  DATE '2026-04-01'
GROUP BY c.region
ORDER BY revenue DESC;
\end{lstlisting}

\begin{block}{需要解释的不是语法，而是口径}
对象范围、连接键、时间窗口、分组维度、度量公式、排序含义。
\end{block}
\end{frame}

\begin{frame}{SQL 的逻辑处理顺序}
\[
\texttt{FROM/JOIN}
\rightarrow
\texttt{WHERE}
\rightarrow
\texttt{GROUP BY}
\rightarrow
\texttt{HAVING}
\rightarrow
\texttt{SELECT}
\rightarrow
\texttt{ORDER BY}.
\]

\begin{table}
\small
\centering
\begin{tabular}{p{0.24\linewidth}p{0.60\linewidth}}
\toprule
子句 & 管理分析含义\\
\midrule
FROM/JOIN & 确定事实表与维度来源。\\
WHERE & 过滤原始事件或对象。\\
GROUP BY & 声明指标输出粒度。\\
HAVING & 过滤聚合后的分组结果。\\
SELECT & 定义最终字段和指标公式。\\
\bottomrule
\end{tabular}
\end{table}
\end{frame}

\begin{frame}{连接的核心风险：改变事实粒度}
若产品维度在 \texttt{product\_id} 上唯一，则
\[
|\texttt{orders}\Join\texttt{products}|=|\texttt{orders}|.
\]
若维度键重复，连接会复制订单事实。

\begin{table}
\small
\centering
\begin{tabular}{p{0.25\linewidth}p{0.55\linewidth}}
\toprule
风险 & 指标后果\\
\midrule
维度键不唯一 & 销售额、订单数、延期率被放大。\\
外键缺失 & 内连接会丢失无法匹配的事实。\\
粒度混合 & 分母不稳定，部门之间不可比较。\\
\bottomrule
\end{tabular}
\end{table}
\end{frame}

\begin{frame}{KPI 是一组被治理的定义}
\[
M=(O,\mathcal G,\mathcal T,N,D,u,\rho).
\]

\begin{table}
\small
\centering
\begin{tabular}{p{0.14\linewidth}p{0.67\linewidth}}
\toprule
符号 & 含义\\
\midrule
\(O\) & 被度量的对象或事件集合。\\
\(\mathcal G\) & 分组维度，如地区、产品类别、客户分层。\\
\(\mathcal T\) & 时间窗口和时间字段。\\
\(N,D\) & 分子、分母或可加度量。\\
\(u,\rho\) & 单位与责任归属。\\
\bottomrule
\end{tabular}
\end{table}
\end{frame}

\begin{frame}{比率型指标不能直接平均}
总体毛利率应在目标粒度重新计算：
\[
\frac{\sum_g N_g}{\sum_g D_g}
\ne
\frac{1}{|\mathcal G|}\sum_g\frac{N_g}{D_g}.
\]

\begin{block}{实践规则}
保存销售额、毛利、准时订单数、交付订单数等分子分母；展示百分比时再计算比率。
\end{block}
\end{frame}

\begin{frame}[fragile]{条件聚合：区域级 KPI}
\begin{lstlisting}
SELECT
  region,
  COUNT(*) AS delivered_orders,
  SUM(revenue) AS revenue,
  SUM(gross_profit) AS gross_profit,
  SUM(gross_profit) / NULLIF(SUM(revenue), 0) AS gross_margin,
  SUM(on_time_flag) * 1.0 / NULLIF(COUNT(*), 0) AS on_time_rate,
  AVG(delay_days) AS average_delay_days
FROM order_enriched
GROUP BY region;
\end{lstlisting}

\begin{alertblock}{安全除法}
分母为零不是普通技术错误，而是需要解释的业务状态。
\end{alertblock}
\end{frame}

\begin{frame}[fragile]{窗口函数：保留行身份的聚合}
\begin{lstlisting}
SELECT
  customer_id,
  order_id,
  order_date,
  revenue,
  SUM(revenue) OVER (
    PARTITION BY customer_id
    ORDER BY order_date
    ROWS BETWEEN 2 PRECEDING AND CURRENT ROW
  ) AS recent_revenue
FROM order_enriched;
\end{lstlisting}

\begin{block}{典型用途}
排名、累计、移动窗口、前后状态比较、客户订单序列。
\end{block}
\end{frame}

\begin{frame}{从 SQL 脚本到指标层}
\begin{table}
\small
\centering
\begin{tabular}{p{0.26\linewidth}p{0.55\linewidth}}
\toprule
字段 & 说明\\
\midrule
name / label & 稳定名称与展示名称。\\
grain & 指标对应的事实粒度。\\
numerator / denominator & 比率型指标的分子分母。\\
dimensions & 可下钻维度。\\
time dimension & 时间窗口字段。\\
owner & 业务解释责任人。\\
\bottomrule
\end{tabular}
\end{table}
\end{frame}

\begin{frame}{第3周实验交付}
\begin{enumerate}
  \item 加载 `data/sample` 中的订单、客户、产品和库存数据；
  \item 构造订单级收入、毛利、准时标记和延期天数；
  \item 按区域、客户分层、产品类别和渠道计算 KPI；
  \item 输出至少一个窗口函数结果；
  \item 为一个指标补全指标字典记录；
  \item 用 Python 复核 SQL 结果。
\end{enumerate}
\end{frame}

\begin{frame}{本讲小结}
\begin{enumerate}
  \item SQL 是管理指标口径的可执行表达；
  \item 连接前必须检查事实粒度是否保持稳定；
  \item KPI 需要对象范围、维度、时间、分子、分母和责任；
  \item 比率型指标应在目标粒度重新计算；
  \item 窗口函数适合表达经营节奏；
  \item 下一讲进入数据质量与指标治理。
\end{enumerate}
\end{frame}

\end{document}

